\chapter{Cursor shapes and visibility}
\label{ch:cursor}

\begin{aside}
The cursor is a single cell of state that controls how the user
feels about input. Slow it down and the app feels sluggish.
Blink it where it shouldn't be and the app feels nervous. Hide
it when not needed and the screen feels calm.
\end{aside}

\section{Six shapes}

DECSCUSR (DEC Set Cursor Style, \code{ESC [ N q}) picks the
shape:

\begin{itemize}[leftmargin=*]
  \item \code{0} or \code{1}: blinking block (default).
  \item \code{2}: steady block.
  \item \code{3}: blinking underline.
  \item \code{4}: steady underline.
  \item \code{5}: blinking bar (I-beam).
  \item \code{6}: steady bar.
\end{itemize}

\maya{} exposes a \code{set\_cursor\_shape(...)} call. By
convention, vim-like modal apps use:

\begin{itemize}[leftmargin=*]
  \item Steady block for normal mode.
  \item Steady bar for insert mode.
  \item Steady underline for visual/select mode.
\end{itemize}

\section{Visibility}

\code{ESC [ ? 25 l} hides; \code{ESC [ ? 25 h} shows. \maya{}'s
default is hidden, because most TUI screens are not text inputs
and the cursor would be a distracting artefact.

A text input element shows the cursor when focused and hides
it otherwise. \maya{} tracks ``the focused input wants the
cursor at $(x, y)$'' and emits the position at end of frame
along with the shape and visibility state.

\section{Position}

\code{ESC [ y ; x H} moves the cursor (1-based row and column).
\maya{} emits this at the end of every frame to position the
cursor where the focused input wants it, then shows it.

If no input has focus, no position is emitted and the cursor
stays hidden.

\section{Blink}

The terminal does the blinking; the application has no
control over the rate. xterm's default is 500 ms half-period.

Some users find blinking distracting. Steady shapes are
available for them. \maya{} does not enforce a blink/steady
choice; it follows the application's request.

\begin{funfact}
The IBM 2741 terminal (1965) had a mechanical cursor: a
solenoid pushed a small block of metal up to mark the
typing position. Blinking is a much later invention; early
CRTs simply held the position.
\end{funfact}

\section{Cursor and the diff}

The diff in Chapter~\ref{ch:diffloop} emits cursor moves as part
of run-to-run navigation. After the diff, the cursor is at the
end of the last run. If the focused input wants it elsewhere,
one more move is emitted.

For typical UIs this final move is 8 bytes of ANSI. Cheap.

\section{Restoring on exit}

On exit, \maya{} emits \code{ESC [ 0 q} (reset cursor shape)
and \code{ESC [ ? 25 h} (show). These are part of the
shutdown ANSI bundle along with alt-screen exit, mouse-mode
off, and bracketed-paste off.

If the program crashes without running the shutdown, the
cursor may be left in a strange shape. \maya{} installs a
\code{std::atexit} hook to be safe; the hook is idempotent.

\begin{warning}
A crashed app with cursor hidden is hard for users to
recover from. Always install the show-cursor hook even if
the rest of cleanup is debatable.
\end{warning}

\section{Multiple cursors}

The terminal has one cursor. To simulate multiple cursors
(for column editing, for instance) you must draw them yourself
as styled cells (reverse video, bold, or a coloured block) and
let the real cursor sit invisible.

\maya{} provides \code{caret} as an element that draws a fake
cursor at a logical position. Use it for editor-like UIs that
need more than one insertion point.

\section{The blink-off frame problem}

Some terminals send a special ``blink off'' redraw every half-
second to clear the cursor cell. If your diff has been
disturbed in the cursor cell, the blink-off can leave a stale
glyph briefly.

DEC 2026 synchronised output fixes this for capable terminals.
For others, accept the half-second artefact.

\section{Cursor and selection}

When showing a text selection (highlighted range), the cursor
sits at one end. The selection is drawn as styled cells
(reverse video, typically); the cursor is the active end.

\maya{}'s text-input widget tracks a head (cursor) and an
anchor (other end). Drawing the selection is a styled fill of
the range; positioning the cursor is one move at end of frame.

\section{Cross-terminal behaviour}

\begin{itemize}[leftmargin=*]
  \item \textbf{xterm:} all six DECSCUSR shapes work.
  \item \textbf{Apple Terminal:} only block; bar and underline
    do nothing.
  \item \textbf{Windows Terminal:} all six; modern.
  \item \textbf{tmux:} passes through; depends on the outer
    terminal.
  \item \textbf{screen:} mostly passes through; some older
    versions drop DECSCUSR.
\end{itemize}

\maya{} sends the request unconditionally. Terminals that
don't understand it ignore the sequence.

\section{Recap}

\begin{recap}
\begin{itemize}[leftmargin=*]
  \item Six DECSCUSR shapes plus a visibility toggle.
  \item Hidden by default; shown by focused inputs.
  \item Position is part of end-of-frame ANSI.
  \item Restore shape and visibility on exit, always.
  \item Multiple cursors are simulated, not real.
\end{itemize}
\end{recap}

\section*{Workshop}

\begin{enumerate}[leftmargin=*]
  \item Build a modal app that switches cursor shape on
    mode change. Confirm in xterm and in your daily
    terminal.
  \item Make the cursor blink rate user-configurable... and
    discover you cannot, then write that down as a
    constraint.
  \item Install a crash test (assert false) and confirm the
    cursor is restored to a block-and-visible state on the
    way out.
\end{enumerate}
